% CUMCM (全国大学生数学建模竞赛) electronic paper skeleton for a Mathodology workspace.
% Rules: 全国大学生数学建模竞赛论文格式规范（2026年修订稿）, checked 2026-11-05:
%   https://www.mcm.edu.cn/html_cn/node/4cd596519c9eb9fbd866398f6df0caa3.html
% Not an official template. Re-read the rules of the year you enter; see
% ../../../../references/contest-formats.md. Compile with XeLaTeX.
%
% Electronic version: the first page is the abstract page (no 承诺书, no 编号专用页).
% No table of contents. Pages are numbered from the abstract page, centred in the footer.
% No name, school or region anywhere. Body at most 30 pages; appendix unlimited.
\documentclass[UTF8,12pt,a4paper]{ctexart}
\usepackage[margin=2.5cm]{geometry}
\usepackage{amsmath,amssymb,booktabs,graphicx}
\usepackage[hidelinks]{hyperref}
\usepackage{gbt7714}
\bibliographystyle{gbt7714-numerical}
\graphicspath{{../figures/}}
\input{results}
\ctexset{section/format=\centering\zihao{4}\heiti}
\pagestyle{plain}

\newcommand{\PaperTitle}{论文标题}

\begin{document}

% ---- 摘要专用页：第 1 页，原则上不超过一页 ----
\begin{center}
  {\zihao{3}\heiti \PaperTitle}\par\medskip
  {\zihao{4}\heiti 摘\quad 要}
\end{center}
% 写清每个问题的方法与结论，数值用 \res{q1.key} 引用。
针对问题一，……

\medskip
\noindent{\heiti 关键词：}关键词一；关键词二；关键词三
\clearpage

% ---- 正文（不超过 30 页，不要目录）----
\section{问题重述}
\section{模型假设}
\section{符号说明}
\section{问题一的模型建立与求解}
% \begin{table}[htbp]\centering\caption{...}\label{tab:q1}\input{tables/q1-compare}\end{table}
\section{模型的评价}

% \bibliography{references}

% ---- 附录：支撑材料文件清单 + 全部可运行源程序 ----
\clearpage
\appendix
\section*{附录}
支撑材料文件清单（无支撑材料时写“本论文没有支撑材料”）。

源程序（未用到程序时写“本论文没有用到程序”）。
\end{document}
